\documentclass[10pt,a4paper]{amsart}
\usepackage{amsmath,amssymb,amsthm,graphicx}
% Retained minimal preamble for the Stacks Project chapter more-algebra.tex (More on Algebra).
% The chapter file pinned-original-source.tex carries only \input{preamble} on its first line, so no
% preamble text lives inside the chapter itself.  The Stacks Project uses one single shared preamble for
% every chapter; the snapshot stages it identically as sources/stacks-chow/preamble.tex and
% sources/stacks-etale-cohomology/preamble.tex.  The block below is copied mechanically, line for line,
% from sources/stacks-chow/preamble.tex lines 14-18 (the Xy-pic setup) and 154-217 (the native theorem
% declarations together with the mathematical macros).  Nothing was added, removed or redefined; the
% wrapper class and the amsmath/amssymb/amsthm/graphicx base packages are supplied by the assembler.
% Chapter: more-algebra.tex -- More on Algebra.
% For commutative diagrams we use Xy-pic
\usepackage[all]{xy}

% We use 2cell for 2-commutative diagrams.
\xyoption{2cell}

% Theorem environments.
%
\theoremstyle{plain}
\newtheorem{theorem}[subsection]{Theorem}
\newtheorem{proposition}[subsection]{Proposition}
\newtheorem{lemma}[subsection]{Lemma}

\theoremstyle{definition}
\newtheorem{definition}[subsection]{Definition}
\newtheorem{example}[subsection]{Example}
\newtheorem{exercise}[subsection]{Exercise}
\newtheorem{situation}[subsection]{Situation}

\theoremstyle{remark}
\newtheorem{remark}[subsection]{Remark}
\newtheorem{remarks}[subsection]{Remarks}

\numberwithin{equation}{subsection}

% Macros
%
\def\lim{\mathop{\mathrm{lim}}\nolimits}
\def\colim{\mathop{\mathrm{colim}}\nolimits}
\def\Spec{\mathop{\mathrm{Spec}}}
\def\Hom{\mathop{\mathrm{Hom}}\nolimits}
\def\Ext{\mathop{\mathrm{Ext}}\nolimits}
\def\SheafHom{\mathop{\mathcal{H}\!\mathit{om}}\nolimits}
\def\SheafExt{\mathop{\mathcal{E}\!\mathit{xt}}\nolimits}
\def\Sch{\mathit{Sch}}
\def\Mor{\mathop{\mathrm{Mor}}\nolimits}
\def\Ob{\mathop{\mathrm{Ob}}\nolimits}
\def\Sh{\mathop{\mathit{Sh}}\nolimits}
\def\NL{\mathop{N\!L}\nolimits}
\def\CH{\mathop{\mathrm{CH}}\nolimits}
\def\proetale{{pro\text{-}\acute{e}tale}}
\def\etale{{\acute{e}tale}}
\def\QCoh{\mathit{QCoh}}
\def\Ker{\mathop{\mathrm{Ker}}}
\def\Im{\mathop{\mathrm{Im}}}
\def\Coker{\mathop{\mathrm{Coker}}}
\def\Coim{\mathop{\mathrm{Coim}}}

% Boxtimes
%
\DeclareMathSymbol{\boxtimes}{\mathbin}{AMSa}{"02}

%
% Macros for moduli stacks/spaces
%
\def\QCohstack{\mathcal{QC}\!\mathit{oh}}
\def\Cohstack{\mathcal{C}\!\mathit{oh}}
\def\Spacesstack{\mathcal{S}\!\mathit{paces}}
\def\Quotfunctor{\mathrm{Quot}}
\def\Hilbfunctor{\mathrm{Hilb}}
\def\Curvesstack{\mathcal{C}\!\mathit{urves}}
\def\Polarizedstack{\mathcal{P}\!\mathit{olarized}}
\def\Complexesstack{\mathcal{C}\!\mathit{omplexes}}
% \Pic is the operator that assigns to X its picard group, usage \Pic(X)
% \Picardstack_{X/B} denotes the Picard stack of X over B
% \Picardfunctor_{X/B} denotes the Picard functor of X over B
\def\Pic{\mathop{\mathrm{Pic}}\nolimits}
\def\Picardstack{\mathcal{P}\!\mathit{ic}}
\def\Picardfunctor{\mathrm{Pic}}
\def\Deformationcategory{\mathcal{D}\!\mathit{ef}}

% --- appended 2026-10-12 (Chief): the project own definitions of the reference/slogan/history
% display environments, copied mechanically from the shared Stacks preamble lines 8-12
% (\usepackage{verbatim} plus the three \newenvironment lines, which the chapter preamble keeps
% verbatim from the project preamble). Nothing added, removed or redefined: these environments
% carry no mathematical content and the project itself typesets them as comments. Without these
% three lines a counted statement whose own text uses them fails with Environment slogan undefined.
\usepackage{verbatim}
\newenvironment{reference}{\comment}{\endcomment}
\newenvironment{slogan}{\comment}{\endcomment}
\newenvironment{history}{\comment}{\endcomment}
\usepackage{cleveref}
\crefname{originalequation}{original equation}{original equations}
\Crefname{originalequation}{Original equation}{Original equations}
\makeatletter
\newlabel{equation-complex-bocksteins}{{27686}{source}{}{sourcepointer-1}{}}
\newlabel{equation-complex-bocksteins@cref}{{[originalequation][0][]27686}{[1][0][]source}{}{}{}}
\makeatother
\crefname{originalremark}{original remark}{original remarks}
\Crefname{originalremark}{Original remark}{Original remarks}
\makeatletter
\newlabel{remark-eta-BZ}{{27579}{source}{}{sourcepointer-2}{}}
\newlabel{remark-eta-BZ@cref}{{[originalremark][0][]27579}{[1][0][]source}{}{}{}}
\makeatother
\title{The eta transform reduced modulo the element is quasi-isomorphic to the Bockstein complex}
\author{The Stacks Project Authors}
\date{stacks-project:more-algebra.tex:local-snapshot-20261012; GNU FDL 1.2 or later}
\begin{document}
\maketitle
\setcounter{subsection}{1}
\noindent\textit{Editorial note.} All delivered mathematical text is the author's own, copied from the pinned Stacks Project chapter file more-algebra.tex; no statement or proof step was rewritten, reordered or removed.\par\medskip
\noindent\textit{Editorial note.} The chapter file pinned-original-source.tex carries only \textbackslash{}input\{preamble\} on its first line, so no preamble text lives inside the chapter. The retained file sources/stacks-more-algebra/retained-preamble.tex is a mechanical, line-for-line copy of the project's single shared preamble as staged in the snapshot; nothing was added, removed or redefined and no package was added.\par\medskip
\noindent\textit{Editorial note.} The source URL is the source-level Stacks Project URL, as required when more than one card is produced from one source; the per-claim tag is carried by claim\_key.\par\medskip
\noindent\textit{Editorial note.} The counted statement span refers to the Bockstein complex by its numbered equation label, delivered verbatim. The author notes inside the counted proof span that all tensor products are over A; that sentence is reproduced. The two references inside the counted proof span are recorded as non-load-bearing navigation pointers.\par\medskip
\noindent\textit{Allowed standard prerequisites (editorial).}\begin{itemize}
\item The eta transform eta\_f M\textasciicircum{}bullet of a complex of f-torsion free modules and its grading. used at proof lines 27753-27754. Supplies the source complex and the description of its degree i part as f\textasciicircum{}i M\textasciicircum{}i.
\item The Bockstein complex H\textasciicircum{}bullet(M\textasciicircum{}bullet/f) with its differentials beta. named in the statement at line 27748 and used at proof lines 27768-27769. Is the target of the canonical map and the source of the compatibility computation.
\item Vanishing of the kernel of the comparison maps (eta\_f M)\textasciicircum{}i tensor A/fA -> Z\textasciicircum{}i/B\textasciicircum{}i and acyclicity. invoked at proof lines 27774-27779. Is what upgrades the map of complexes to a quasi-isomorphism.
\end{itemize}\medskip
\ifdefined\hypertarget\hypertarget{sourcepointer-1}{}\fi
\noindent\textit{Source reference.} Original equation 27686: sources/stacks-more-algebra/pinned-original-source.tex, source lines 27686--27686. Same-chapter navigation pointer; the statement it names is standard for the section and the argument that consumes it stays inside the counted proof span.\par\medskip
\ifdefined\hypertarget\hypertarget{sourcepointer-2}{}\fi
\noindent\textit{Source reference.} Original remark 27579: sources/stacks-more-algebra/pinned-original-source.tex, source lines 27579--27579. Same-chapter navigation pointer; the statement it names is standard for the section and the argument that consumes it stays inside the counted proof span.\par\medskip
\begin{lemma}
\label{lemma-eta-second-property}
Let $A$ be a ring and let $f \in A$ be a nonzerodivisor. Let $M^\bullet$ be
a complex of $f$-torsion free $A$-modules.
There is a canonical map of complexes
$$
\eta_fM^\bullet \otimes_A A/fA
\longrightarrow
H^\bullet(M^\bullet/f)
$$
which is a quasi-isomorphism where the right hand side is the complex
(\ref{equation-complex-bocksteins}).
\end{lemma}

\begin{proof}
Let $x \in (\eta_fM)^i$. Then $x = f^ix' \in f^iM$ and
$d^i(x) = f^{i + 1}y \in f^{i + 1}M^{i + 1}$. Thus $d^i$ maps $x' \otimes f^i$
to zero in $M^{i + 1} \otimes f^iA/f^{i + 1}A$. All tensor products
are over $A$ in this proof.
Hence we may map $x$ to the class of $x' \otimes f^i$ in
$H^i(M^\bullet \otimes f^iA/f^{i + 1}A)$. It is clear that this
rule defines a map
$$
(\eta_fM)^i \otimes A/fA
\longrightarrow
H^i(M^\bullet \otimes f^iA/f^{i + 1}A)
$$
of $A/fA$-modules. Observe that in the situation above, we may view
$x' \otimes f^i$ as an element of $M^i \otimes f^iA/f^{i + 2}A$
with differential $d^i(x' \otimes f^i) = y \otimes f^{i + 1}$.
By the construction of $\beta$ above we find that
$\beta(x' \otimes f^i) = y \otimes f^{i + 1}$ and we conclude that
our maps are compatible with differentials, i.e., we have a map
of complexes.

\medskip\noindent
To finish the proof, we observe that the construction given
in the previous paragraph agrees with the maps
$(\eta_fM)^i \otimes A/fA \to Z^i/B^i$ discussed in
Remark \ref{remark-eta-BZ}.
Since we have seen that the kernel of these maps is
an acyclic subcomplex of $\eta_fM^\bullet \otimes A/fA$, the lemma is proved.
\end{proof}

\end{document}
